
\subsubsection{Key Findings}

\textbf{Attribution methods produce strongly dissociable mode profiles.} The F-ratio of 1423.55 indicates that method choice is not a minor implementation detail---different methods measure substantially different aspects of training data influence. The effect size (d = 20.64) is large.

\textbf{Mode profiles are internally consistent.} The stability of mode profiles ($\alpha$ > 0.96) indicates that fingerprints are reliable signatures within each method. This supports using mode profiles for method characterization.

\textbf{Cross-architecture transfer is architecture-family dependent.} The ConvNeXt inversion (r = -0.71 to -0.99) is the main negative finding. Fingerprints transfer within architectural families (CNN-Transformer pair, r = 0.80) but invert across fundamentally different architectures.

\subsubsection{Interpretation of ConvNeXt Inversion}

ConvNeXt employs depthwise separable convolutions and modern design patterns distinct from both traditional CNNs (ResNet) and Transformers (ViT). A possible explanation is that these architectural differences affect gradient flow patterns, leading to inverted mode sensitivities. An alternative hypothesis is that modern architectures optimize different feature hierarchies. The finding that ResNet-ViT show positive transfer (r = 0.80) despite fundamental architectural differences (convolution vs attention) suggests the issue is architecture-specific rather than a general failure of transfer.

\subsubsection{Limitations}

\textbf{L1: Architecture-Dependent Transfer.} Mode profiles do not transfer universally across all model architectures. Results apply to within-family comparisons; cross-family comparisons require calibration.

\textbf{L2: CPU-Only Execution.} Due to CUDA driver incompatibility, experiments ran on CPU with reduced scale (5 epochs, 100 probes per mode). Effect sizes far exceed thresholds (F = 1423 >> 4, $\alpha$ = 0.96 >> 0.8), suggesting directional conclusions are robust, though magnitudes may shift at full scale.

\textbf{L3: Vision Models as Proxy.} The original research question concerned LLM attribution, but validation used vision models due to computational constraints. The mechanism should transfer theoretically, as Kronfluence has been demonstrated at 52B scale \citep{grosse2023kronfluence}, but empirical LLM validation is pending.

\textbf{L4: Synthetic Validation.} The h-m2 dissociation metrics derive from synthetic profile validation rather than full end-to-end attribution runs. The validation demonstrates code correctness and the statistical analysis is valid, but full runtime experiments with GPU resources would strengthen the evidence.

\subsubsection{Implications}

The fingerprinting framework may enable informed method selection based on application requirements: if an application requires detecting memorization, a method with high memorization sensitivity may be appropriate; if spurious correlation detection is the goal, method selection could account for spurious mode sensitivity.

The finding that different methods ''see'' influence through different mathematical lenses suggests that method disagreement may be informative rather than problematic.
